\begin{figure}[htbp]
\centering
\begin{tikzpicture}[x=1cm,y=1cm,font=\small,
  ray/.style={-{Latex[length=2mm]},very thick,draw=NimbusBlue}]
  \draw[fill=NimbusLight,draw=NimbusBlue] (0,0) rectangle (4,3);
  \draw[draw=NimbusBlue] (0,1)--(4,1) (0,2)--(4,2);
  \node[above] at (2,3.65) {Column interfaces carry flux};
  \draw[ray] (.8,3.5)--(.8,2.55);
  \node[left] at (.7,3.25) {$F^\downarrow$};
  \draw[ray] (3.2,2.55)--(3.2,3.5);
  \node[right] at (3.3,3.25) {$F^\uparrow$};
  \node at (2,2.5) {layer $k+1$};
  \node at (2,1.5) {layer $k$: $\rho_d,h,T$};
  \node at (2,.5) {layer $k-1$};
  \draw[very thick,draw=NimbusBlue] (-.2,0)--(4.2,0);
  \node[below=5pt] at (2,0) {surface: $T_s,\epsilon,\alpha$};
  \draw[ray] (4.15,1.5)--(5.15,1.5);
  \node[draw=NimbusBlue,fill=NimbusLight,rounded corners,align=center,
        text width=6cm,minimum height=2cm,anchor=west] at (5.3,1.5)
    {Interface difference\\[4pt]
     $\rho_d c_p h\dot T=F_{k-1/2}-F_{k+1/2}$\\[5pt]
     publish $\dot\theta$ and cell power together};
\end{tikzpicture}
\caption{Upward-positive net flux is $F=F^\uparrow-F^\downarrow$.
The layer gains energy when more net power enters its lower face than leaves
its upper face. Summing layers cancels interior exchanges. Optical closure,
surface boundary conditions and physical layer mass remain separate assumptions.}
\end{figure}
